\chapter{Methodology}
In this chapter an 


\section{System Architecture}
This section will detail the overall architecture of the federated learning system implemented in this study.
It will describe the components of the system, how they interact with each other, and the rationale behind architectural choices.
The focus will be on illustrating the structure that supports federated learning operations, including client-server communication,
data distribution, and model aggregation.

\subsection{Component Overview}
\subsection{Interaction Flow}

%%%%%%%%%%%%%%%%%%%%%%%%%%%%%%%%%%%%%%%%%%%%%%%%%%%%%%%%%%%%%%%%%%%%%%%%%%%%%%%
\section{Dataset Preprocessing \& Management}
In this section the different datasets used and the criteria for their selection.
Emphasis will be placed on the preprocessing steps undertaken to prepare the data
for analysis, ensuring the data's compatibility with the requirements of federated
learning.

\subsection{Dataset Sources and Selection Criteria}
\subsection{Preprocessing Techniques}

%%%%%%%%%%%%%%%%%%%%%%%%%%%%%%%%%%%%%%%%%%%%%%%%%%%%%%%%%%%%%%%%%%%%%%%%%%%%%%%
\section{Model Development and Training}
for different datasets different model designes are implemented that need to be described.
Moreover the custom techniques of the training process has to be exaplained.

\subsection{Model Architecture}
\subsection{Training and Optimization}

%%%%%%%%%%%%%%%%%%%%%%%%%%%%%%%%%%%%%%%%%%%%%%%%%%%%%%%%%%%%%%%%%%%%%%%%%%%%%%%
\section{Federated Learning Implementation}
Here, we detail the implementation of the federated learning framework, focusing on the client-server architecture,
the distribution of the model across clients, and the aggregation of model updates. Special attention is given to
modifications or custom implementations designed to enhance the efficiency of federated learning in this research,
as multiple adjustmnents needed to be conducted on the initial strategy

\subsection{Client-Server Setup}
\subsection{Data Partitioning and Distribution}
\subsection{Model Aggregation}

%%%%%%%%%%%%%%%%%%%%%%%%%%%%%%%%%%%%%%%%%%%%%%%%%%%%%%%%%%%%%%%%%%%%%%%%%%%%%%%
\section{Communication Protocols and Data Serialization}
The section discusses the communication protocols established between the clients both the flower and metric server
for efficient data exchange and model synchronization. It covers the use of ProtoBuffers
to optimize network communication, ensuring data integrity and security during transmission.

\subsection{Implementation of Communication Protocols}
\subsection{Data Serialization Technique}

%%%%%%%%%%%%%%%%%%%%%%%%%%%%%%%%%%%%%%%%%%%%%%%%%%%%%%%%%%%%%%%%%%%%%%%%%%%%%%%
\section{Network Simulation Integration}
In this section, we delve into the use of the NS3 simulator to replicate real-world network conditions and their impact
on federated learning. It details the simulation setup, the parameters adjusted for different network conditions, and
how these simulations contribute to understanding federated learning's behavior in varied networking environments.

\subsection{Simulation Tool Configuration}
\subsection{Simulating Network Conditions}
\subsection{Integration with Federated Learning}

%%%%%%%%%%%%%%%%%%%%%%%%%%%%%%%%%%%%%%%%%%%%%%%%%%%%%%%%%%%%%%%%%%%%%%%%%%%%%%%
\section{Evaluation Metrics and Performance Analysis}
This section outlines the metrics used to evaluate the performance of the federated learning models in conjuction with
the network simulation stats.

\subsection{Model Evaluation Metrics}
\subsection{Network Simulation Metrics}
\subsection{Performance Analysis}
%%%%%%%%%%%%%%%%%%%%%%%%%%%%%%%%%%%%%%%%%%%%%%%%%%%%%%%%%%%%%%%%%%%%%%%%%%%%%%%
